\section{Feature dictionaries and data availability}
\label{sec:si:features}

The two samples are described by different feature sets drawn from different pulls, so each has
its own dictionary. Table~\ref{tab:archivefeatures} lists the \ArcFeatures{} features of the
primary sample with the FBref table each column was read from, and Table~\ref{tab:features}
lists the \NFeatures{} features of the reduced sample, which combines what survived the
withdrawal with an independent expected goals source. In both, counting statistics are
expressed per ninety minutes played and rate features are used as published rather than divided
a second time.

Every feature in the reduced set was confirmed to hold non-null values in both of its seasons
before being admitted, which is the check that distinguishes a live column from one whose values
were withdrawn. Table~\ref{tab:availability} reports that audit. The primary set is audited
differently, by per-column and per-season coverage on load, because its failure mode is a block
missing from one season rather than a column emptied across all of them.

\begin{table}[t]
\centering
\small
\caption{The primary feature set, with the FBref table each column was read from. Counting statistics are per ninety minutes played; the seven rate features are used as published. Five defensive counts are possession adjusted, as described in Section~\ref{sec:methods:padj}.}
\label{tab:archivefeatures}
\begin{tabular}{llp{7.0cm}}
\toprule
Feature & FBref table & Definition \\
\midrule
np xg p90 & standard & Non-penalty expected goals (per 90) \\
shots p90 & shooting & Shots taken, excluding penalties (per 90) \\
goals non penalty p90 & standard & Goals excluding penalties (per 90) \\
xag p90 & standard & Expected assisted goals, the expected goals of shots a pass created (per 90) \\
key passes p90 & passing & Passes leading directly to a shot (per 90) \\
sca p90 & gca & Shot creating actions, the two offensive actions before a shot (per 90) \\
gca p90 & gca & Goal creating actions, the two offensive actions before a goal (per 90) \\
passes final third p90 & passing & Completed passes into the final third (per 90) \\
passes penalty area p90 & passing & Completed passes into the penalty area (per 90) \\
crosses p90 & misc & Crosses attempted (per 90) \\
progressive passes p90 & passing & Completed passes moving the ball substantially toward goal (per 90) \\
progressive carries p90 & possession & Carries moving the ball substantially toward goal (per 90) \\
progressive receptions p90 & possession & Passes received in a progressive position (per 90) \\
carries final third p90 & possession & Carries into the final third (per 90) \\
touches def pen p90 & possession & Touches in the defending team's own penalty area (per 90) \\
touches def third p90 & possession & Touches in the defensive third (per 90) \\
touches mid third p90 & possession & Touches in the middle third (per 90) \\
touches att third p90 & possession & Touches in the attacking third (per 90) \\
touches att pen p90 & possession & Touches in the opposition penalty area (per 90) \\
take ons p90 & possession & Take-ons attempted (per 90) \\
tackles padj p90 & defense & Tackles, possession adjusted (per 90) \\
interceptions padj p90 & misc & Interceptions, possession adjusted (per 90) \\
blocks padj p90 & defense & Blocks of a pass or shot, possession adjusted (per 90) \\
clearances padj p90 & defense & Clearances, possession adjusted (per 90) \\
ball recoveries p90 & misc & Loose balls recovered (per 90) \\
fouls committed padj p90 & misc & Fouls committed, possession adjusted (per 90) \\
shot accuracy pct & shooting & Shots on target as a share of shots \\
take on success pct & possession & Take-ons completed as a share of those attempted \\
tackle win pct & defense & Duels won as a share of those contested \\
aerials won pct & misc & Aerial duels won as a share of those contested \\
pass cmp short pct & passing & Short passes completed as a share of those attempted \\
pass cmp medium pct & passing & Medium passes completed as a share of those attempted \\
pass cmp long pct & passing & Long passes completed as a share of those attempted \\
\bottomrule
\end{tabular}
\end{table}


\begin{table}[t]
\centering
\small
\caption{The outfield feature set. Counting statistics are per ninety minutes. The two rate features are used as published.}
\label{tab:features}
\begin{tabular}{llp{7.2cm}}
\toprule
Feature & Source & Definition \\
\midrule
np\_xg\_p90 & Understat & Non-penalty expected goals \\
xa\_p90 & Understat & Expected assists \\
key\_passes\_p90 & Understat & Passes leading directly to a shot \\
xg\_chain\_p90 & Understat & Expected goals of possessions the player was involved in \\
xg\_buildup\_p90 & Understat & As xGChain, excluding shots and key passes \\
shots\_p90 & FBref & Shots taken, excluding penalties \\
shots\_on\_target\_p90 & FBref & Shots on target \\
goals\_non\_penalty\_p90 & FBref & Goals excluding penalties \\
assists\_p90 & FBref & Assists \\
crosses\_p90 & FBref & Crosses attempted \\
interceptions\_padj\_p90 & FBref & Interceptions, possession adjusted \\
tackles\_won\_padj\_p90 & FBref & Tackles won, possession adjusted \\
fouls\_committed\_padj\_p90 & FBref & Fouls committed, possession adjusted \\
fouls\_drawn\_p90 & FBref & Fouls drawn \\
offsides\_p90 & FBref & Times caught offside \\
cards\_yellow\_p90 & FBref & Yellow cards \\
shot\_accuracy\_pct & FBref & Shots on target as a share of shots, a rate \\
goals\_per\_shot & FBref & Goals per shot, a rate \\
\bottomrule
\end{tabular}
\end{table}


\begin{table}[t]
\centering
\small
\caption{Populated columns in each FBref player table. The tables are still served with complete headers, so the empty counts are the number of columns that resolve correctly and contain no values.}
\label{tab:availability}
\begin{tabular}{lrrrr}
\toprule
 & \multicolumn{2}{c}{2024-25} & \multicolumn{2}{c}{2025-26} \\
\cmidrule(lr){2-3}\cmidrule(lr){4-5}
Table & Columns & Empty & Columns & Empty \\
\midrule
standard & 17 & 0 & 17 & 0 \\
shooting & 11 & 0 & 11 & 0 \\
passing & 21 & 19 & 21 & 19 \\
passing types & 16 & 14 & 16 & 14 \\
goal shot creation & 17 & 16 & 17 & 16 \\
defense & 17 & 14 & 17 & 14 \\
possession & 21 & 20 & 21 & 20 \\
playing time & 17 & 0 & 17 & 0 \\
misc & 13 & 2 & 13 & 2 \\
keeper & 19 & 0 & 19 & 0 \\
keeper adv & 26 & 23 & 26 & 23 \\
\bottomrule
\end{tabular}
\end{table}


\section{Possession elasticities}
\label{sec:si:elasticity}

Before any structure is estimated, the defensive counts have to be separated from the
opportunity to record them. Team defensive volume does rise as opponent possession rises, so the
confound is real, and the conventional correction, in the proportional form that has succeeded
the original of \citet{knutson2014possession}, assumes the relationship is proportional: a side
with twice the ball is taken to concede half the defensive events.

It is much weaker than that for \NElastExcludeOne{} of the five statistics, and it is not
distinguishable from
proportional for the fifth. Fitted on the primary sample over \ArcPadjTeamSeasons{}
team-seasons, with \ElastNBootstrap{} bootstrap resamples giving each estimate an interval, the
elasticities are \ArcElastTackles{} for tackles (\ElastCITacklesLow{} to \ElastCITacklesHigh{}),
\ArcElastInterceptions{} for interceptions (\ElastCIInterceptionsLow{} to
\ElastCIInterceptionsHigh{}), \ArcElastBlocks{} for blocks (\ElastCIBlocksLow{} to
\ElastCIBlocksHigh{}) and \ArcElastFouls{} for fouls committed (\ElastCIFoulsLow{} to
\ElastCIFoulsHigh{}). All four intervals exclude one, and no bootstrap resample of any of them
reached it. Clearances are different: the estimate is \ArcElastClearances{}, the interval runs
from \ElastCIClearancesLow{} to \ElastCIClearancesHigh{}, \ElastShareOneClearances{} percent of
resamples land at or above one, and the bootstrap probability that the true elasticity is at
least one is \ElastPOneClearances{}. For clearances the conventional exponent cannot be
rejected, and the paper does not claim otherwise.

That pattern fits the football. A clearance is the most purely reactive action in the set, made
when the ball arrives rather than sought, so it is the statistic for which the exposure story
should hold best, and it is the one where proportionality survives. A tackle is the most
volitional, and it is furthest from proportional. The spread from \ArcElastMin{} to
\ArcElastMax{} is itself the point: a single exponent applied to every defensive statistic is
wrong for four of them.

The reduced sample supplies the check the primary sample cannot, because it retains the
team-level possession series needed to measure what is left after correcting. Across its
\PadjTeamSeasons{} team-seasons, team possession ranges from \PadjPossessionMin{} to
\PadjPossessionMax{} percent, and its fitted elasticities agree with the primary sample in
falling well below one, at \ElasticityInterceptions{} for interceptions,
\ElasticityTackles{} for tackles won and \ElasticityFouls{} for fouls committed.
Table~\ref{tab:padj} shows what each choice of exponent leaves behind. The three raw
correlations with team possession are \PadjCorrRawInterceptions{}, \PadjCorrRawTackles{} and
\PadjCorrRawFouls{}, all negative, since a player whose team keeps the ball records fewer
defensive actions. Applying the conventional exponent overshoots every one of them, giving
\PadjCorrUnitInterceptions{}, \PadjCorrUnitTackles{} and \PadjCorrUnitFouls{}, reversing sign
in all three cases and, for fouls committed, ending further from zero than where it started.
The correction passes through the confound and out the other side. Adjusting with the measured
elasticity instead leaves \PadjCorrAdjInterceptions{}, \PadjCorrAdjTackles{} and
\PadjCorrAdjFouls{}, three to four times closer to zero than the conventional exponent achieves.

Two qualifications belong with this. The residual-dependence check runs on one league, whose
possession shares are relatively compressed, so the exact residuals should not be transported to
a competition with a wider spread without refitting. Refitting the elasticities within each
league tests transport directly, and the answer is mixed. For tackles, interceptions, fouls and
clearances every league's interval overlaps the pooled one. Blocks do not transport: the league
estimates run from \ElastLeagueMinBlocks{} in \ElastLeagueMinNameBlocks{} to
\ElastLeagueMaxBlocks{} in \ElastLeagueMaxNameBlocks{}, and the intervals for both of those
leagues miss the pooled interval entirely. Clearances are also the statistic where leagues
disagree most, with \ElastLeaguesAboveOneClearances{} of the five league point estimates at or
above one, which is a second reason not to claim sub-proportionality for them.

The constant elasticity form is itself an assumption, and for \NElastQuadMatters{} of the five
statistics it fails. Adding a quadratic term in log opponent possession improves the fit for
tackles, blocks and fouls, with a coefficient whose interval excludes zero in each case, and in
each the response flattens toward zero for the sides that see least of the ball. Fouls committed
bend most: the implied elasticity is \ElastAtLowFouls{} at \ElastPossLow{} percent opponent
possession and \ElastAtHighFouls{} at \ElastPossHigh{} percent. The exponents are therefore
averages over the possession range rather than constants. Figure~\ref{fig:sielasticity}
gathers the pooled estimates, their intervals and the league refits, and places both families
of statistic either side of the proportional line.

\begin{figure}[!htbp]
  \centering
  \includegraphics[width=\textwidth]{elasticity_uncertainty.pdf}\par\bigskip
  \includegraphics[width=\textwidth]{possession_elasticities.pdf}
  \caption{Possession elasticities. Top, the five defensive elasticities with bootstrap
  intervals from \ElastNBootstrap{} resamples of \ArcPadjTeamSeasons{} team-seasons, against the
  conventional value of one, with each league's own refit shown as a smaller marker. Bottom,
  both families of statistic on the reduced sample: defensive volume responds to opponent
  possession less than proportionately, while attacking output responds to a team's own
  possession more than proportionately.}
  \label{fig:sielasticity}
\end{figure}

\begin{table}[t]
\centering
\small
\caption{Possession adjustment. The elasticity is fitted from team totals. The final three columns give the correlation between each per-90 rate and team possession: unadjusted, adjusted with the fitted elasticity, and adjusted with the conventional exponent of one. A value near zero is the goal.}
\label{tab:padj}
\begin{tabular}{lrrrr}
\toprule
 & & \multicolumn{3}{c}{Correlation with team possession} \\
\cmidrule(lr){3-5}
Statistic & Elasticity & Unadjusted & Adjusted & Exponent one \\
\midrule
Interceptions & 0.525 & -0.122 & -0.022 & 0.070 \\
Tackles won & 0.432 & -0.111 & -0.018 & 0.105 \\
Fouls committed & 0.283 & -0.102 & -0.037 & 0.129 \\
\bottomrule
\end{tabular}
\end{table}


\subsection{Why the correction is one sided}
\label{sec:si:symmetric}

Correcting only the defensive statistics invites an obvious objection. Attacking output also
depends on possession, so the feature set is still team dependent and the correction looks like
half a correction. The objection deserves a test rather than an assertion. The test runs on the
reduced sample, which is the only one of the two carrying an independent expected-goals source
with involvement measures, and its conclusion concerns the shape of the two relationships rather
than their exact magnitudes.

The two sides of the ball do not behave alike. Fitted the same way, against a team's own
possession rather than the opponent's, the attacking elasticities are \BetaShots{} for shots,
\BetaNPXG{} for non-penalty expected goals, \BetaGoals{} for non-penalty goals, \BetaAssists{}
for assists, \BetaXGChain{} for xGChain and \BetaXGBuildup{} for xGBuildup. Every one except
crossing, at \BetaCrosses{}, exceeds one, while \NElastExcludeOne{} of the five defensive
elasticities fall
well below it and the fifth cannot be distinguished from it.

That asymmetry is the whole argument. For defending, opponent possession is an exposure
denominator in the plain sense that a tackle cannot occur while your own side holds the ball,
and the sub-proportional response says exposure is what it measures. For attacking, an
elasticity above one says the opposite: ten percent more of the ball buys appreciably more than
ten percent more output. Returns to possession are increasing, which is a statement about
football rather than about measurement, so dividing it out would remove signal along with any
bias. The two involvement measures make the point most starkly, since both exceed two.

The counterfactual was nonetheless computed. Applying the fitted exponents to the attacking
block as well moves the correlation between club centroids and final league position from
\SymCorrBaseline{} to \SymCorrAdjusted{}, so a symmetric correction does remove a real part of
the team relationship rather than nothing at all. It does not remove all of it, the partition
moves at an adjusted Rand index of \SymARI{} against the one reported here, and the silhouette
drifts from \SymSilhouetteBase{} to \SymSilhouetteAdj{}. The choice is a judgement rather than a
fact, and it is recorded as one: we treat possession driven productivity as something a role
profile should reflect, and a reader who takes the opposite view has the numbers to see what
that would cost.

\section{Feature-family ablation in full}
\label{sec:si:ablation}

Table~\ref{tab:ablation} gives every cell of the leave-one-family-out and keep-one-family-only
ablation summarised in Section~\ref{M-sec:results:replication}: for each feature set and each
scope, the number of clusters the declared rule selects, the standard score of the observed
two-cluster silhouette against the clusterless null, and the adjusted Rand index between the
ablated partition and the full one on the same players. The last of those is what distinguishes
a surviving count of two clusters from the same two modes surviving.

\begin{table}[t]
\centering
\small
\caption{Standard score of the observed two-cluster silhouette against the clusterless Gaussian null, by scope, under every feature-family ablation. A dagger marks a cell where the pre-declared rule chose a number of clusters other than two.}
\label{tab:ablation}
\begin{tabular}{lrrrr}
\toprule
Feature set & All & DF & MF & FW \\
\midrule
\multicolumn{5}{l}{\emph{Full feature set}} \\
all 33 features & 18.0 & 24.2 & 15.0 & 13.2 \\
\midrule
\multicolumn{5}{l}{\emph{Leave one family out}} \\
without shooting & 17.7 & 21.3 & 12.3 & 18.0 \\
without creation & 21.2 & 21.1 & 14.5 & 15.1 \\
without progression & 19.0 & 19.7 & 13.2 & 19.8 \\
without territory & 21.6 & 22.6 & 16.0 & 15.6 \\
without defending & 19.7 & 29.1 & 16.0 & 17.2 \\
without passing & 19.1 & 17.5 & 10.2 & 20.1 \\
\midrule
\multicolumn{5}{l}{\emph{One family alone}} \\
shooting only & 29.3 & 19.8 & 16.1 & 2.9 \\
creation only & 12.8 & 22.6 & 11.7 & 18.4 \\
progression only & 8.3 & 15.6 & 9.0 & 12.9 \\
territory only & 11.4 & 26.6 & 16.2 & 20.2 \\
defending only & 19.0 & 6.9 & 10.7 & 8.5 \\
passing only & 6.6 & 11.4 & 5.3 & 1.4 \\
\bottomrule
\end{tabular}
\end{table}


\section{Nesting: resampling players rather than player-seasons}
\label{sec:si:nesting}

The unit of analysis is the player-season, and \ArcEligible{} of them come from
\NestPlayersAll{} players, so an average player supplies \NestPerPlayerAll{} rows. A bootstrap
that resamples rows therefore treats correlated observations as independent, and the stability
it reports is optimistic by an unknown amount. This section measures the amount.

The identical bootstrap is run twice on the same partition with the same number of resamples,
\NestNBoot{}. The first draws player-seasons with replacement, which is what the rest of the
paper reports and what the genre does. The second draws whole players with replacement and takes
every season a drawn player contributed, so that rows sharing a player are drawn or left out
together. The adjusted Rand index is computed the same way in both, on the distinct rows the
draw contains.

The difference is small. On the outfield pool the row bootstrap returns \NestRowAll{} and the
player block bootstrap \NestBlockAll{}. Within position groups the block index is
\NestBlockDF{} for defenders, \NestBlockMF{} for midfielders and \NestBlockFW{} for forwards,
against \NestRowDF{}, \NestRowMF{} and \NestRowFW{} for the row bootstrap. The largest fall
anywhere is \NestLargestFall{}, for forwards, which is the smallest scope and the one the paper
already reports as the least well supported.

Two qualifications belong with the comparison. A player block bootstrap draws fewer distinct
rows than a row bootstrap of the same nominal size, because one drawn player contributes several
correlated rows at once, so a small fall is expected from the reduced effective sample alone and
is not by itself evidence that the row bootstrap was misleading. The check also applies to the
observed data only: the clusterless datasets of Section~\ref{M-sec:methods:null} are independent
rows by construction, with no player to block on, so the comparison between observed and null
stability is unaffected.

What the check establishes is narrow and worth having. The stability reported for this partition
does not depend on treating player-seasons as exchangeable, so the argument that stability fails
to distinguish real structure from a partitioned continuum is not an artefact of the resampling
unit either.

\section{Minutes threshold sensitivity}
\label{sec:si:sensitivity}

Both samples exclude outfield players below \MinMinutes{} minutes. That threshold balances two
errors, since a low threshold admits players whose per-90 rates are dominated by sampling noise
while a high threshold discards genuine squad roles such as the impact substitute.
Table~\ref{tab:sensitivity} repeats the entire selection and stability procedure at 270,
\MinMinutes{} and 900 minutes, treating each threshold as an independent pool in which
imputation and standardisation are recomputed and the declared rule is reapplied from scratch.

The threshold matters very little. The rule selects two clusters at every threshold, the
silhouette moves only from \SensSilhouetteLow{} to \SensSilhouetteHigh{}, and bootstrap
stability is between \SensARILow{} and \SensARIHigh{} throughout. The partition itself barely
moves: measured on the players common to both pools, the adjusted Rand index against the
partition at the threshold actually used is \SensAgreeLow{} at 270 minutes and \SensAgreeHigh{}
at 900. Figure~\ref{fig:sisensitivity} plots the same quantities.

\begin{table}[t]
\centering
\small
\caption{Sensitivity of the analysis to the minutes threshold in the primary season.}
\label{tab:sensitivity}
\begin{tabular}{rrrrr}
\toprule
Minutes & Eligible & $k$ & Silhouette & Bootstrap ARI \\
\midrule
270 & 403 & n/a & n/a & n/a \\
450 & 366 & n/a & n/a & n/a \\
900 & 296 & n/a & n/a & n/a \\
\bottomrule
\end{tabular}
\end{table}


\begin{figure}[!htbp]
  \centering
  \includegraphics[width=\textwidth]{minutes_sensitivity.pdf}
  \caption{Eligible sample, silhouette and bootstrap stability across the candidate minutes
  thresholds. The dashed line marks the threshold used in the main analysis.}
  \label{fig:sisensitivity}
\end{figure}

\section{Selection criteria and embedding choices}
\label{sec:si:selection}

Table~\ref{tab:kselection} gives the silhouette coefficient, Davies-Bouldin index,
Calinski-Harabasz index and gap statistic for every candidate number of clusters on the reduced
sample, for the global analysis and within each position group, with the row the declared rule
selects marked in each scope. The corresponding evidence for the primary sample is
Figure~\ref{M-fig:structurenull}, which carries more information than a criteria table would,
because the question the primary sample answers is not which candidate scores best but whether
any of them beats what a continuum produces.

Table~\ref{tab:umapgrid} reports trustworthiness for all twelve configurations of the neighbour
count and minimum distance, on the reduced sample, which is the only one for which a manifold
embedding is reported. Selection was by trustworthiness alone, decided before any embedding was
inspected visually, which matters because a manifold method will produce a visually appealing
picture at almost any setting.

\begin{longtable}{llrrrrr}
\caption{Cluster selection criteria for every candidate number of clusters, by scope. The consensus rule is applied to these values.}\\
\label{tab:kselection}\\
\toprule
Scope & $k$ & Silhouette & Davies-Bouldin & Calinski-Harabasz & Gap \\
\midrule
\endfirsthead
\toprule
Scope & $k$ & Silhouette & Davies-Bouldin & Calinski-Harabasz & Gap \\
\midrule
\endhead
global$\ast$ & 2 & 0.287 & 1.451 & 145.6 & 1.216 \\
global & 3 & 0.215 & 1.695 & 101.9 & 1.283 \\
global & 4 & 0.154 & 1.954 & 86.1 & 1.339 \\
global & 5 & 0.145 & 1.861 & 75.4 & 1.381 \\
global & 6 & 0.144 & 1.797 & 68.3 & 1.415 \\
global & 7 & 0.135 & 1.905 & 61.3 & 1.438 \\
global & 8 & 0.118 & 1.881 & 56.1 & 1.458 \\
global & 9 & 0.123 & 1.840 & 53.5 & 1.477 \\
global & 10 & 0.123 & 1.805 & 49.8 & 1.484 \\
global & 11 & 0.128 & 1.657 & 47.6 & 1.506 \\
global & 12 & 0.121 & 1.911 & 44.6 & 1.517 \\
DF$\ast$ & 2 & 0.204 & 2.037 & 25.8 & 0.943 \\
DF & 3 & 0.151 & 1.919 & 23.5 & 1.011 \\
DF & 4 & 0.142 & 1.983 & 19.7 & 1.036 \\
DF & 5 & 0.120 & 1.961 & 18.0 & 1.059 \\
DF & 6 & 0.117 & 1.795 & 16.9 & 1.083 \\
DF & 7 & 0.106 & 1.898 & 15.6 & 1.109 \\
DF & 8 & 0.120 & 1.792 & 14.7 & 1.136 \\
DF & 9 & 0.120 & 1.754 & 14.1 & 1.154 \\
DF & 10 & 0.112 & 1.636 & 13.5 & 1.159 \\
DF & 11 & 0.098 & 1.636 & 12.8 & 1.162 \\
DF & 12 & 0.105 & 1.627 & 12.3 & 1.196 \\
MF$\ast$ & 2 & 0.243 & 1.543 & 69.4 & 1.035 \\
MF & 3 & 0.147 & 1.969 & 49.0 & 1.091 \\
MF & 4 & 0.145 & 1.895 & 40.4 & 1.130 \\
MF & 5 & 0.145 & 1.873 & 34.7 & 1.153 \\
MF & 6 & 0.132 & 1.939 & 30.3 & 1.167 \\
MF & 7 & 0.122 & 2.054 & 27.4 & 1.199 \\
MF & 8 & 0.107 & 2.030 & 24.9 & 1.202 \\
MF & 9 & 0.103 & 1.923 & 23.8 & 1.212 \\
MF & 10 & 0.106 & 1.983 & 21.9 & 1.220 \\
MF & 11 & 0.094 & 1.944 & 20.4 & 1.227 \\
MF & 12 & 0.097 & 1.835 & 19.5 & 1.226 \\
FW$\ast$ & 2 & 0.187 & 1.803 & 12.4 & 0.764 \\
FW & 3 & 0.175 & 1.576 & 10.8 & 0.837 \\
FW & 4 & 0.176 & 1.599 & 10.2 & 0.908 \\
FW & 5 & 0.196 & 1.181 & 9.3 & 0.909 \\
FW & 6 & 0.124 & 1.490 & 8.5 & 0.944 \\
FW & 7 & 0.132 & 1.350 & 7.8 & 1.017 \\
FW & 8 & 0.101 & 1.389 & 7.7 & 1.021 \\
FW & 9 & 0.127 & 1.229 & 7.7 & 1.070 \\
FW & 10 & 0.125 & 1.350 & 7.8 & 1.051 \\
FW & 11 & 0.128 & 1.199 & 7.1 & 1.064 \\
FW & 12 & 0.116 & 1.126 & 6.8 & 1.096 \\
\bottomrule
\end{longtable}


\begin{table}[t]
\centering
\small
\caption{UMAP hyperparameter grid, ranked by trustworthiness. Selection was by trustworthiness alone and was decided before any embedding was inspected.}
\label{tab:umapgrid}
\begin{tabular}{rrr}
\toprule
Neighbours & Minimum distance & Trustworthiness \\
\midrule
10 & 0 & 0.9337 \\
15 & 0 & 0.9322 \\
15 & 0.10 & 0.9318 \\
30 & 0 & 0.9312 \\
10 & 0.10 & 0.9309 \\
50 & 0 & 0.9273 \\
30 & 0.10 & 0.9242 \\
10 & 0.50 & 0.9218 \\
50 & 0.10 & 0.9179 \\
15 & 0.50 & 0.9134 \\
50 & 0.50 & 0.9091 \\
30 & 0.50 & 0.9084 \\
\bottomrule
\end{tabular}
\end{table}


\section{Fuller views of the main-text figures}
\label{sec:si:profiles}

Figure~\ref{fig:sigeometry} gives the three panels of Figure~\ref{M-fig:geometry} at full size and
with the detail a six-panel composite cannot carry: the principal component loadings across all
five leading components rather than the first two, the mode against position contingency
normalised by row as well as by column, and the confusion matrices for both classifiers rather
than only the boosted one. Figure~\ref{fig:siprofiles} gives the centroid deviation of each archetype from its position
group mean across the entire feature set, rather than only the ten features with the largest
separation shown in Figure~\ref{M-fig:modes}. Figure~\ref{fig:siradars} gives the same centroids
as radars on fixed axes within each position group. Within a group the two profiles are mirror
images by construction, since the features are standardised within group before clustering, so
the useful reading is of shape rather than of extent.

Figure~\ref{fig:sishap} gives the full feature attribution for the position classifier. The
direction of each attribution adds a detail that magnitude alone hides. Ranking the correlation
between a feature's value and the attribution it makes to a class separates a class identified
by what its members do from one identified by what they do not, and the three classes differ on
exactly that. \ShapPresenceDF{} of the five strongest features for defenders are positively
signed, so defenders are recognised by the presence of what they do. Forwards and midfielders
are recognised as much by the absence of defensive volume as by the presence of attacking
output, with only \ShapPresenceFW{} and \ShapPresenceMF{} of their five positive respectively.
A class defined largely by negations is a class occupying one end of a scale rather than a
category with a signature of its own. Figure~\ref{fig:sishapclass} splits the same attribution
by class.

\begin{figure}[!p]
  \centering
  \includegraphics[width=\textwidth]{archive_pca.pdf}
  \caption{Principal components of the \ArcFeatures{} archive features over \ArcEligible{}
  player-seasons, with the scree curve and the loadings of the leading components at full size. Figure~\ref{fig:sicontingency} gives the
  contingency and Figure~\ref{fig:siconfusion} the confusion matrices.}
  \label{fig:sigeometry}
\end{figure}

\begin{figure}[!p]
  \centering
  \includegraphics[width=\textwidth]{archive_cluster_vs_position.pdf}
  \caption{The pooled two-mode partition against listed position, with the heat maps normalised
  by row so that each gives the position composition of one mode. The main text reads the same
  table by column, which is where the \ArcSplitDF{} and \ArcSplitFW{} percent figures come
  from.}
  \label{fig:sicontingency}
\end{figure}

\begin{figure}[!htbp]
  \centering
  \includegraphics[width=\textwidth]{archive_confusion.pdf}
  \caption{Out-of-fold confusion for both position classifiers on \SupN{} player-seasons. The
  near equality of the linear and the boosted model is itself evidence that little interaction
  structure remains for the trees to find.}
  \label{fig:siconfusion}
\end{figure}

\begin{figure}[!p]
  \centering
  \includegraphics[width=\textwidth,height=0.86\textheight,keepaspectratio]{archive_archetype_profiles.pdf}
  \caption{Centroid deviation of each archetype from its position group mean across the full
  \ArcFeatures{} feature set, primary sample. Within a group the two profiles are mirror images
  by construction.}
  \label{fig:siprofiles}
\end{figure}

\begin{figure}[!p]
  \centering
  \includegraphics[width=\textwidth]{archive_radar_panel_DF.pdf}\par\bigskip
  \includegraphics[width=\textwidth]{archive_radar_panel_MF.pdf}
  \caption{Archetype radars on fixed axes, primary sample. Top, defenders. Bottom,
  midfielders. Forwards are in Figure~\ref{fig:siradarsfw}.}
  \label{fig:siradars}
\end{figure}

\begin{figure}[!htbp]
  \centering
  \includegraphics[width=\textwidth]{archive_radar_panel_FW.pdf}
  \caption{Archetype radars for forwards on the primary sample, on the same fixed axes as
  Figure~\ref{fig:siradars}.}
  \label{fig:siradarsfw}
\end{figure}

\begin{figure}[!p]
  \centering
  \includegraphics[width=\textwidth,height=0.88\textheight,keepaspectratio]{archive_shap_beeswarm.pdf}
  \caption{SHAP values for the gradient boosted position classifier on the primary sample,
  showing for each player the contribution of each feature toward the position group the player
  is actually listed at. Territorial features dominate technical ones throughout.}
  \label{fig:sishap}
\end{figure}

\begin{figure}[!p]
  \centering
  \includegraphics[width=\textwidth,height=0.88\textheight,keepaspectratio]{archive_shap_by_class.pdf}
  \caption{Mean absolute attribution per feature, split by position group, primary sample.
  Touches in the defensive third dominate for all three, and the profile of what matters differs
  sharply between them: defenders are separated largely by defensive volume and crossing,
  forwards by where they receive the ball. The panel shows magnitude only, and the direction
  each feature pushes is reported in the text.}
  \label{fig:sishapclass}
\end{figure}

\section{Players the description fits worst}
\label{sec:si:outliers}

A two-mode description of a continuum will fit some players badly, and which players those are
is a test of whether the description is measuring football or noise. Isolation forests and
Mahalanobis distance under a shrinkage covariance were fitted separately within each position
group, and the two detectors agree on \ArcOutAgreement{} percent of the players they flag.

The \ArcOutN{} strongest outliers are \ArcOutDistinct{} distinct players, headed by
\ArcOutTop{}. Two features of the list stand out. \ArcOutCompound{} of the \ArcOutN{} carry a
compound listed position such as midfielder and forward together, against a much smaller share
in the pool, so the players the partition handles worst are the ones the listing system also
declines to place. The detector is also not simply finding stars: it finds attacking full backs
whose output profile resembles a winger's, which is the same boundary the defender split
describes and a genuine tactical category rather than an artefact.

One caveat belongs with the list rather than after it. \ArcOutShortMinutes{} of the \ArcOutN{}
played under nine hundred minutes in the season concerned, so their per-90 rates carry more
sampling noise than a full season's would, and extreme rates are easier to reach on a small
denominator. The list should be read as evidence about the kind of player the two-mode
description struggles with, not as a ranking in which the fifteenth entry means what the first
does. Figure~\ref{fig:sioutliers} shows both detectors against each other.

\begin{figure}[!htbp]
  \centering
  \includegraphics[width=\textwidth]{archive_outliers.pdf}
  \caption{Isolation forest score against Mahalanobis distance within position group, primary
  sample. The two detectors are computed independently and rank the same players highly, and the
  flagged points are dominated by players carrying compound listed positions.}
  \label{fig:sioutliers}
\end{figure}

\section{Statistics that change meaning by position}
\label{sec:si:divergence}

Correlations were computed within each position group and compared across groups, and the ten
pairs whose correlation varies most between groups are reported in Table~\ref{tab:divergence}
and shown in Figure~\ref{fig:sidivergence}. The largest spread between position groups reaches
\CorrSpreadMax{}, against a median spread across all pairs of \CorrSpreadMedian{}.

Every one of the ten most divergent pairs changes sign between position groups. The relationship
between shots on target and crosses is positive for midfielders and negative for forwards, which
is to say that a midfielder who crosses often is generally a more active attacker while a
forward who crosses often is generally not the player getting on the end of chances. The
relationship between key passes and tackles won is positive for defenders and negative for
midfielders, distinguishing the full back whose defensive work and creative output rise together
from the midfield division of labour between a destroyer and a creator.

The practical implication is that a single pooled correlation matrix, of the kind that appears
routinely in public football analysis, averages over these sign reversals and reports a number
that describes no actual population of players. It also supplies the justification for the
position stratified analysis used throughout this paper, since a global embedding fitted across
all outfield players is fitted across relationships that point in opposite directions. The full
correlation matrix the comparison is drawn from is Figure~\ref{fig:siheatmap}.

\begin{table}[t]
\centering
\small
\caption{The ten feature pairs whose Spearman correlation differs most across position groups. Every pair changes sign between at least two groups.}
\label{tab:divergence}
\begin{tabular}{llrrrr}
\toprule
Feature A & Feature B & DF & MF & FW & Spread \\
\midrule
Shots on target & Crosses & -0.01 & 0.43 & -0.44 & 0.87 \\
Crosses & Shot accuracy \% & -0.22 & 0.26 & -0.57 & 0.82 \\
npxG & Crosses & -0.15 & 0.36 & -0.43 & 0.79 \\
Key passes & Tackles won & 0.45 & -0.34 & 0.23 & 0.79 \\
Crosses & Interceptions & 0.20 & -0.50 & 0.26 & 0.76 \\
Crosses & Tackles won & 0.50 & -0.25 & 0.47 & 0.75 \\
xA & Shots on target & 0.17 & 0.56 & -0.19 & 0.75 \\
Shots & Interceptions & 0.19 & -0.52 & -0.05 & 0.70 \\
xA & Tackles won & 0.39 & -0.32 & 0.11 & 0.70 \\
xA & Interceptions & 0.18 & -0.51 & 0.11 & 0.70 \\
\bottomrule
\end{tabular}
\end{table}


\begin{figure}[!htbp]
  \centering
  \includegraphics[width=\textwidth]{corr_divergence_top10.pdf}
  \caption{The ten feature pairs whose within-group correlation differs most across position
  groups, in the earlier of the two reduced-sample seasons. Every pair changes sign between at
  least two groups.}
  \label{fig:sidivergence}
\end{figure}

\begin{figure}[!p]
  \centering
  \includegraphics[width=\textwidth]{corr_heatmap_all.pdf}
  \caption{The full hierarchically clustered Spearman correlation matrix over the reduced
  outfield feature set, from which the comparison in Figure~\ref{fig:sidivergence} is drawn.}
  \label{fig:siheatmap}
\end{figure}

\section{Analyses on the reduced sample}
\label{sec:si:reduced}

The analyses collected here run on the reduced sample described in
Section~\ref{M-sec:methods:data}: the seasons ending 2025 and 2026, pulled live after the
withdrawal, with roughly half the features of the primary sample.

One point about scope needs stating plainly, because it applies to every figure in this section.
Where a figure here shows Premier League players only, that is not a sampling choice made for
convenience. It is the coverage of the live post-withdrawal pull that this robustness case is
built on, and the reduced sample exists to test whether the main result survives a much poorer
description of the same game rather than to characterise European football. The paper's claims
rest on the primary sample, which covers all \ArcNLeaguesWord{} leagues and all \ArcNSeasonsWord{}
seasons, and every one of them is refitted independently in each league and each season in
Section~\ref{M-sec:results:replication}. Section~\ref{sec:si:leagues} below additionally repeats
the reduced-sample pipeline across all five leagues, so the reduced sample's own generality is
tested rather than assumed. The within-season drift analysis of Section~\ref{sec:si:drift} is
narrower still, covering one club, and the reason is stated there.

\subsection{Cluster selection and stability}
\label{sec:si:clustering}

Applying the declared rule to the four selection criteria gives $k=\KGlobal{}$ globally and
within each position group, matching the primary sample. Across \BootstrapN{} bootstrap
resamples the mean adjusted Rand index against the reference partition is \BootARIGlobal{}
globally and \BootARIMF{} for midfielders, both describing a partition that reappears almost
unchanged under resampling. For defenders the mean falls to \BootARIDF{} with a standard
deviation of \BootARISDDF{}, and for forwards to \BootARIFW{} with a standard deviation of
\BootARISDFW{}. The minimum adjusted Rand index observed across forward resamples is
\BootARIMinFW{}, which is to say that in the least favourable resample the partition recovered
is no better than chance. Forwards number only \NFWPrimary{} players above the minutes threshold
here, so this is partly a sample size problem, and it is a real limit on what the reduced sample
can claim. The global partition scores an adjusted Rand index of \ARIvsPosition{} against the
position group with a cluster purity of \PurityVsPosition{}, so the clusters are not a
rediscovery of the positions the source already records. Figure~\ref{fig:sikselection} gives
the criterion curves, Figure~\ref{fig:siconsensus} the co-assignment structure and
Figure~\ref{fig:siclustervalidation} the correspondence with position.

\subsection{Descriptive structure and embeddings}
\label{sec:si:descriptive}

After filtering, \NEligiblePrimary{} outfield players are eligible in the earlier of the two
reduced-sample seasons, comprising \NDFPrimary{} defenders, \NMFPrimary{} midfielders and
\NFWPrimary{} forwards. Summary statistics by position group appear in Table~\ref{tab:summary}.
The feature set is strongly non-normal: all \NFeaturesRejectNormal{} of the \NFeatures{}
features reject the Shapiro-Wilk null at the five percent level, mean absolute skewness is
\SkewMeanAbs{}, and Mardia's test rejects multivariate normality decisively on both components.
Principal component analysis requires no distributional assumption to be computed, so this does
not invalidate the decomposition, but the components should be read as directions of variance
rather than as latent factors, and it justifies the use of Spearman rather than Pearson
correlation throughout.

Principal component analysis, UMAP and t-SNE were fitted to the same standardised features. The
selected UMAP configuration used \UMAPBestNeighbors{} neighbours at a minimum distance of
\UMAPBestMinDist{}. Trustworthiness is \TrustPCA{} for the linear projection, \TrustUMAP{} for
UMAP and \TrustTSNE{} for t-SNE, while continuity is close to identical across the three.
Nearest neighbour position purity at ten neighbours is \PurityPCA{} for the linear projection
and \PurityUMAP{} for UMAP, so roughly two thirds to three quarters of a player's closest
statistical neighbours share his listed position group. That figure is informative in both
directions: it confirms the surviving features carry real positional signal, and it establishes
that a substantial minority of every player's nearest neighbours are listed at a different
position, which is the space in which data-driven roles can differ from listed ones. All three
embeddings display a continuous gradient from defenders through midfielders to forwards rather
than separated groups, and that gradient is the honest description of the geometry.
Figure~\ref{fig:sidescriptive} shows the distributions and the three embeddings and
Figure~\ref{fig:siredpca} the principal component structure, while
Figures~\ref{fig:siredarchetypes} to \ref{fig:siredradarsfw} give the archetypes this sample
returns.

\begin{table}[t]
\centering
\small
\caption{Headline per-90 statistics by position group in the primary season, reported as mean with standard deviation in parentheses.}
\label{tab:summary}
\begin{tabular}{lrrr}
\toprule
Statistic & DF & MF & FW \\
\midrule
np\_xg\_p90 & 0.05 (0.04) & 0.16 (0.12) & 0.41 (0.17) \\
xa\_p90 & 0.06 (0.06) & 0.15 (0.10) & 0.16 (0.11) \\
key\_passes\_p90 & 0.51 (0.44) & 1.26 (0.63) & 1.05 (0.54) \\
xg\_chain\_p90 & 0.35 (0.16) & 0.50 (0.20) & 0.63 (0.20) \\
xg\_buildup\_p90 & 0.30 (0.14) & 0.28 (0.13) & 0.18 (0.09) \\
shots\_p90 & 0.54 (0.33) & 1.54 (0.85) & 2.53 (0.69) \\
shot\_accuracy\_pct & 29.22 (20.60) & 31.91 (13.77) & 40.30 (9.65) \\
crosses\_p90 & 1.22 (1.62) & 2.26 (1.75) & 1.40 (1.62) \\
interceptions\_p90 & 0.95 (0.37) & 0.74 (0.44) & 0.24 (0.16) \\
tackles\_won\_p90 & 1.06 (0.48) & 1.17 (0.55) & 0.55 (0.33) \\
fouls\_committed\_p90 & 0.87 (0.38) & 1.32 (0.55) & 1.27 (0.62) \\
offsides\_p90 & 0.06 (0.07) & 0.16 (0.17) & 0.53 (0.33) \\
\bottomrule
\end{tabular}
\end{table}


\begin{figure}[!htbp]
  \centering
  \includegraphics[width=\textwidth]{distributions_by_position.pdf}\par\bigskip
  \includegraphics[width=\textwidth]{embedding_comparison.pdf}
  \caption{Reduced sample. Top, distributions of four representative per-90 statistics across
  position groups, where separation is clear for shot volume and defensive actions and much
  weaker for creation. Bottom, two-dimensional embeddings of the same feature set coloured by
  listed position group; the defender to forward gradient is continuous in all three, and
  t-SNE is shown for visual contrast only and is never used for clustering.}
  \label{fig:sidescriptive}
\end{figure}

\begin{figure}[!p]
  \centering
  \includegraphics[width=\textwidth]{archetype_profiles.pdf}
  \caption{Centroid deviation of each archetype from its position group mean, reduced sample.}
  \label{fig:siredarchetypes}
\end{figure}

\begin{figure}[!p]
  \centering
  \includegraphics[width=\textwidth]{radar_panel_DF.pdf}\par\bigskip
  \includegraphics[width=\textwidth]{radar_panel_MF.pdf}
  \caption{Archetype radars on the reduced sample. Top, defenders. Bottom, midfielders. The two
  sets of radars in this document are drawn on different feature counts, so their axes are not
  directly comparable and the useful comparison is of shape rather than extent.}
  \label{fig:siredradars}
\end{figure}

\begin{figure}[!htbp]
  \centering
  \includegraphics[width=\textwidth]{radar_panel_FW.pdf}
  \caption{Archetype radars for forwards on the reduced sample, on the same axes as
  Figure~\ref{fig:siredradars}.}
  \label{fig:siredradarsfw}
\end{figure}

\begin{figure}[!htbp]
  \centering
  \includegraphics[width=\textwidth]{pca_scree.pdf}\par\bigskip
  \includegraphics[width=\textwidth]{pca_biplot.pdf}
  \caption{Principal component structure of the reduced outfield feature set.
  \PCNinetyComponents{} components are required to reach ninety percent of the variance, so the
  feature set is not close to low rank, and the first three components together account for
  \PCCumThree{}. Top, scree and cumulative variance. Bottom, the loadings biplot.}
  \label{fig:siredpca}
\end{figure}

\begin{figure}[!htbp]
  \centering
  \includegraphics[width=\textwidth]{k_selection_curves.pdf}
  \caption{Cluster selection criteria across candidate cluster counts, reduced sample. The
  consensus rule was fixed before fitting and selects the same value in every scope.}
  \label{fig:sikselection}
\end{figure}

\begin{figure}[!p]
  \centering
  \includegraphics[width=\textwidth]{consensus_matrix.pdf}
  \caption{Bootstrap co-assignment structure on the reduced sample, one panel per scope, with
  players ordered by cluster.}
  \label{fig:siconsensus}
\end{figure}

\begin{figure}[!htbp]
  \centering
  \includegraphics[width=\textwidth]{cluster_vs_position.pdf}\par\bigskip
  \includegraphics[width=\textwidth]{cluster_scatter_pca.pdf}
  \caption{Reduced sample. Top, the selected partition against listed positions. Bottom, the
  same partition shown in the principal component space.}
  \label{fig:siclustervalidation}
\end{figure}

\subsection{The same pipeline across the big five}
\label{sec:si:leagues}

Repeating the reduced-sample pipeline unchanged on Spain, Germany, Italy and France for the
earlier of the two seasons returns the same answer everywhere. The declared rule selects
\LgKEng{} clusters in all \NLeagues{} leagues. The silhouette at that choice spans \LgSilMin{}
to \LgSilMax{} across the five, bootstrap stability spans \LgBootARIMin{} to \LgBootARIMax{},
and agreement with listed position groups spans \LgARIPosMin{} to \LgARIPosMax{}. Each of
those three ranges is narrow, and the third is uniformly weak: in no league does the partition
recover the positions the source already records. The elasticity result generalises as well,
since all \LgElastCount{} defensive elasticities fitted independently inside each league lie
between \LgElastMin{} and \LgElastMax{}, every one below the value of one that the conventional
possession adjustment assumes. Figure~\ref{fig:sileagues} shows the comparison.

\begin{figure}[!htbp]
  \centering
  \includegraphics[width=\textwidth]{league_comparison.pdf}
  \caption{The reduced-sample pipeline applied unchanged to each of the big five leagues, with
  England, the competition the reduced sample is otherwise drawn from, in the accent colour. The
  silhouette at the selected number of clusters, bootstrap stability and agreement with listed positions are all close across
  competitions, and the declared rule selects the same number of clusters in every one.}
  \label{fig:sileagues}
\end{figure}

\subsection{What a supervised model sees}
\label{sec:si:supervisedreduced}

Predicting the listed position group from the same \NFeatures{} features, gradient boosting
reaches an accuracy of \LGBMAccuracy{} and a macro F1 of \LGBMMacroF{}, and logistic regression
reaches \LogitAccuracy{} and \LogitMacroF{}, against a majority class baseline of
\BaselineAccuracy{}. The features clearly do carry position, so the weak agreement between
clusters and positions reflects the clustering finding a different and equally real axis rather
than an empty feature set.

The attribution carries a finding that follows directly from the withdrawal. With possession and
passing volume gone, positions are identified largely by absence: defenders by low shooting and
low key passes, forwards by low interceptions, and midfielders by an intermediate profile on
both. A classifier of this kind trained before January 2026 would have identified a defender by
what he did with the ball, and it now identifies him by what he does not do. Players the
classifier places in the wrong group overlap with the players the clustering flagged as hybrids,
at \OverlapLGBM{} for gradient boosting against an enrichment of \EnrichmentLGBM{} over
independence, which is corroboration rather than a strong result because the hybrid flag is
itself broad.

The finishing above role regression illustrates a hazard worth recording. Predicting non-penalty
expected goals per ninety from non-shooting behaviour returns a cross validated coefficient of
determination of \FinishingRSq{}. An earlier specification that also included xGChain returned
\FinishingRSqLeaky{}, which looks far better and is in fact circular, since Understat credits
xGChain for possessions the player finished himself, so the difference between xGChain and
xGBuildup contains the target and correlates with it at \FinishingLeakCorr{}. Even at the honest
value roughly half the variance is unexplained, so the residual leaderboard of
Figure~\ref{fig:sifinishing} is an illustration rather than a measurement of finishing ability.
Figure~\ref{fig:sisupervised} shows the confusion and the per-class importance, and
Figure~\ref{fig:sishapreduced} the full attribution.

\begin{figure}[!htbp]
  \centering
  \includegraphics[width=\textwidth]{confusion_matrix.pdf}\par\bigskip
  \includegraphics[width=\textwidth]{shap_by_class.pdf}
  \caption{Supervised validation of the surviving feature set, reduced sample. Top, confusion
  between listed positions. Bottom, mean feature importance per position.}
  \label{fig:sisupervised}
\end{figure}

\begin{figure}[!p]
  \centering
  \includegraphics[width=\textwidth,height=0.88\textheight,keepaspectratio]{shap_beeswarm.pdf}
  \caption{Global SHAP values for the gradient boosted position classifier on the reduced
  sample. Shooting volume and interceptions carry most of the signal, so a player's position is
  identified largely by the actions he does not perform.}
  \label{fig:sishapreduced}
\end{figure}

\subsection{The third dimension and team signatures}
\label{sec:si:threed}

Two components are a convenient way to draw the space and a poor way to describe it. On the
reduced sample the first two carry shooting volume and wide creation and together account for
\ThreeDVarTwo{} of the variance, while the third carries duels and discipline, loading on fouls
committed, cards and tackles won, and is orthogonal to both of the others, so a planar
projection does not compress it, it deletes it. Retaining it raises captured variance to
\ThreeDVarThree{}. Adding the dimension does not change the headline geometry, since the
defender to forward gradient remains continuous rather than resolving into islands.

The third dimension does change the club picture. Club centroids correlate with final league
position at \ThreeDTeamPCOne{} on the first component and \ThreeDTeamPCTwo{} on the second, but
only \ThreeDTeamPCThree{} on the third. Shooting volume and wide creation both track league
position; the physical dimension of duels, fouls and cards essentially does not. A club does not
finish higher by fouling more or less, which needs stating because the duelling axis is a full
component of player variation carrying almost no information about team success. The
three-dimensional view also settles a structural question by measurement rather than assertion:
drawing each archetype centroid as a direction from the origin, the two centroids within a
position group have a cosine of \CosineDF{} for defenders, \CosineMF{} for midfielders and
\CosineFW{} for forwards. They are exactly anti-parallel, which is forced by standardising
within group before a two-way split, so six archetypes occupy three lines through the origin.

Aggregating players into clubs produces the one result in the reduced sample that is both strong
and stable. The minutes-weighted position of a club's outfield players on the first principal
component correlates with final league position at a Spearman coefficient of \TeamPCOneRho{} in
the earlier season and \TeamPCOneRhoRepl{} in the replication season, close to an exact
reproduction. Since the first component is essentially attacking involvement, this says that
clubs whose players collectively shoot and create more finish higher, which is unsurprising in
itself; the near identical coefficient across two independent seasons shows the team level
summary is measuring something stable even where the player level partitions are not.
Hierarchical clustering on archetype composition groups the \NTeams{} clubs into
\NStyleGroups{} style families. Figures~\ref{fig:sithreed} and \ref{fig:siaxes} show the
three-dimensional views, Figure~\ref{fig:siteams} the club signatures, and
Figure~\ref{fig:siteampanels} each club's players against the rest of the league in the shared
embedding, on axes shared across all twenty panels.

\begin{figure}[!p]
  \centering
  \includegraphics[width=\textwidth]{pca3d_positions.pdf}\par\bigskip
  \includegraphics[width=\textwidth]{pca3d_clusters.pdf}
  \caption{The reduced-sample outfield space in three principal components, each drawn from two
  viewpoints so that structure can be told apart from the accident of a single projection. Top,
  players by listed position group, with grey points on the floor supplying a depth cue. Bottom,
  the selected global partition with centroids marked by crosses.}
  \label{fig:sithreed}
\end{figure}

\begin{figure}[!p]
  \centering
  \includegraphics[width=\textwidth]{pca3d_archetype_axes.pdf}\par\bigskip
  \includegraphics[width=\textwidth]{pca3d_teams.pdf}
  \caption{Top, archetype centroid directions drawn as unit vectors, so the figure shows
  orientation rather than magnitude; each position group forms a single line through the origin,
  so six archetypes describe three axes. Bottom, clubs as minutes-weighted centroids in three
  components, shaded by final league position with darker shading for a higher finish.}
  \label{fig:siaxes}
\end{figure}

\begin{figure}[!p]
  \centering
  \includegraphics[width=\textwidth]{team_dendrogram.pdf}\par\bigskip
  \includegraphics[width=\textwidth]{team_composition.pdf}
  \caption{Club signatures on the reduced sample. Top, style families from hierarchical
  clustering on archetype composition. Bottom, archetype composition by club, ordered by final
  league position.}
  \label{fig:siteams}
\end{figure}

\begin{figure}[!p]
  \centering
  \includegraphics[width=\textwidth]{team_panels.pdf}
  \caption{Each club's outfield players highlighted in the shared embedding, with the rest of
  the league in grey. Axes are shared across all twenty panels, so clubs can be compared
  directly, and they differ visibly in how far their defenders sit toward the midfield region.}
  \label{fig:siteampanels}
\end{figure}

\subsection{Similarity, replication and sampling noise}
\label{sec:si:novel}

Nearest neighbour search within position group recovers football sense for players whose roles
are defined by the surviving statistics, and fails informatively for those whose roles are not.
Full backs return full backs and penalty box forwards return penalty box forwards. The
instructive failures are the midfielders: a holding midfielder's nearest neighbours are creative
attacking midfielders, because with no tackles by pitch third, no blocks, no clearances, no
pressures and no passing volume, the only things left that distinguish him are his set piece
deliveries and key passes. The metric finds the players who share that signature and misses the
entire defensive half of his job, which Figure~\ref{fig:sinovel} makes visible.

\NCommonPlayers{} players are eligible in both reduced-sample seasons. The declared rule
independently selects the same number of clusters in every scope in both seasons, so the
structure replicates, while membership replicates only partly: the adjusted Rand index between
seasons is \CrossSeasonARIGlobal{} globally, falling to \CrossSeasonARIDF{} for defenders,
\CrossSeasonARIMF{} for midfielders and \CrossSeasonARIFW{} for forwards, with the transitions
in Figure~\ref{fig:sitransitions}.

Empirical Bayes shrinkage toward the position group mean, with the degree of shrinkage set by
minutes played, gives the sharpest test of whether the partitions describe playing roles or
sampling noise. The global partition survives: only \ShrinkChangedGlobal{} of
\NEligiblePrimary{} players change cluster, the adjusted Rand index between the original and
shrunken partitions is \ShrinkARIGlobal{}, and the rate at which players move falls
monotonically with minutes played, exactly as the noise account predicts. The within group
partitions fare worse and not equally. For defenders the index falls to \ShrinkARIDF{}, a
partition largely dissolved by removing sampling noise, while midfielders retain
\ShrinkARIMF{} and forwards \ShrinkARIFW{}. Figure~\ref{fig:sishrinkageminutes} shows the
change rate against minutes and Figure~\ref{fig:sishrinkage} the effect on every feature.

Read alongside the bootstrap and replication results, the three tests do not deliver a single
uniform verdict on the reduced sample, and reporting them as though they did would overstate the
case. Defenders fail all three there, forwards fail two of the three, and midfielders pass all
three. The conclusion that survives every test on the reduced sample is narrower than a single
summary sentence would suggest: the global division between attacking and defensive involvement
is a real property of players, the midfield split is trustworthy, and the defender split should
not be used to characterise an individual on that feature set. The primary sample, with twice
the features and six times the players, supports all three splits, which is the clearest single
illustration of what the richer description buys.

\begin{figure}[!p]
  \centering
  \includegraphics[width=\textwidth,height=0.88\textheight,keepaspectratio]{similarity_heatmap.pdf}
  \caption{Nearest neighbour similarity structure within position group, reduced sample.}
  \label{fig:sinovel}
\end{figure}

\begin{figure}[!p]
  \centering
  \includegraphics[width=\textwidth]{cross_season_transitions.pdf}
  \caption{Between-season cluster transitions on the reduced sample, for the
  \NCommonPlayers{} players eligible in both seasons.}
  \label{fig:sitransitions}
\end{figure}

\begin{figure}[!htbp]
  \centering
  \includegraphics[width=\textwidth]{shrinkage_by_minutes.pdf}\par\bigskip
  \includegraphics[width=\textwidth]{outliers_scatter.pdf}
  \caption{Reduced sample. Top, reliability and cluster change rate against minutes played,
  which falls monotonically as the noise account predicts. Bottom, the two outlier detectors
  against each other within position group.}
  \label{fig:sishrinkageminutes}
\end{figure}

\begin{figure}[!p]
  \centering
  \includegraphics[width=\textwidth,height=0.88\textheight,keepaspectratio]{shrinkage_effect.pdf}
  \caption{The effect of empirical Bayes shrinkage on every feature of the reduced sample, with
  the shrunken value against the observed one, one panel per feature.}
  \label{fig:sishrinkage}
\end{figure}

\subsection{Role drift within a season}
\label{sec:si:drift}

Every result elsewhere treats a season as one observation per player, which assumes a player's
profile is homogeneous across it. Match reports make that assumption testable, and because the
withdrawal removed the affected columns from the match schema rather than emptying them, what
remains there is fully populated. A full season is 380 reports at seven seconds each, so the
scrape is restricted to one club, \DriftTeam{}, chosen because its eligible players span all six
archetypes and therefore give drift the widest range to move across. \DriftMatches{} matches
were retrieved and \DriftPlayers{} players accumulated at least one complete rolling window of
\DriftWindow{} matches. This is an illustration of a limitation rather than an estimate of its
size, and it should be read as one.

The match feature set has \DriftFeatures{} of the eighteen features used elsewhere, since the
Understat expected goals block has no match equivalent. Rather than assume the coarser space is
a fair stand-in, we refit the season clustering using only those features and compare it to the
partition reported throughout; the two agree at an adjusted Rand index of \DriftReducedARI{}, so
drift measured within the reduced space is interpretable.

Most players do not drift. \DriftSwitchers{} of \DriftPlayers{}, or \DriftSwitchShare{} percent,
change cluster at least once during the season, and the remainder occupy a single cluster
throughout. The players who stay put are the ones whose roles are least ambiguous, holding
midfielders on one side and the recognised attackers on the other, while the movers are
predominantly full backs and wide players whose attacking involvement varies with the fixture.
For four fifths of players a season aggregate describes a single role and the homogeneity
assumption holds; for the remaining fifth it averages over two.
Figure~\ref{fig:sidrift} shows the trajectories.

\begin{figure}[!htbp]
  \centering
  \includegraphics[width=\textwidth]{role_drift.pdf}
  \caption{Rolling cluster assignment across the season for the \DriftTeam{} players who change
  cluster at least once, on a \DriftWindow{} match window. Players not shown occupied one cluster
  for the whole season.}
  \label{fig:sidrift}
\end{figure}

\subsection{Age and archetype}
\label{sec:si:age}

The two archetypes within a position group differ mainly in attacking involvement, and
involvement is plausibly something a career passes through rather than a fixed trait, so the
association with age is a natural test. It is tested within position group, because archetypes
are nested inside positions and a pooled table would largely recover the fact that forwards and
defenders have different age profiles.

Only defenders show a relationship. The crossing and chance creating defenders average
\AgeGapDF{} years younger than the low involvement group, a difference the rank test places at
\AgeKruskalPDF{}, and the share of the crossing archetype falls monotonically from the youngest
band to the oldest. Midfielders show nothing at all, with a gap of \AgeGapMF{} years at
\AgeKruskalPMF{}, and the global partition is unrelated to age at \AgeGlobalP{}. Forwards point
the same way as defenders but cannot support the claim, since six of the ten cells in their
table have expected counts below five.

The defender result also illustrates a point about test choice that is easy to get wrong. A
chi-square on the banded table returns \AgeChiPDF{} and would ordinarily be read as no
association. That test spends its degrees of freedom on any departure from independence, while
the pattern here is monotonic, and a rank correlation against age recovers it at
\AgeTrendRhoDF{} with a p-value of \AgeTrendPDF{}. The banded test is not wrong, it is
underpowered against the alternative that actually holds, and reporting only the chi-square
would have converted a real finding into a null one. Figure~\ref{fig:siage} shows the shares.

\begin{figure}[!htbp]
  \centering
  \includegraphics[width=\textwidth]{age_archetype.pdf}
  \caption{Share of each age band taken by each archetype within position group, reduced sample.
  The defender panel falls monotonically with age, which a chi-square on bands does not detect
  and a rank correlation does.}
  \label{fig:siage}
\end{figure}

\begin{figure}[!p]
  \centering
  \includegraphics[width=\textwidth]{finishing_above_role.pdf}
  \caption{Residuals from the role model on the reduced sample, which should be read as an
  illustration rather than as a measurement of finishing ability, for the reasons given in the
  text.}
  \label{fig:sifinishing}
\end{figure}

\subsection{Goalkeepers on the reduced sample}
\label{sec:si:keepersreduced}

Post-shot expected goals, sweeper actions, distribution and cross handling are all absent from
the reduced sample, so only shot stopping and workload survive. That version builds two
composites instead of five and substitutes the residual from a regression of saves on shots on
target faced for post-shot expected goals. The substitute adjusts for the volume of work a
goalkeeper receives but not for its difficulty, and it is reported as the weaker instrument it
is. For the primary sample, Figure~\ref{fig:sikeepersfull} gives the shot stopping ranking,
Figure~\ref{fig:sikeepercorr} the correlation structure of the goalkeeper block and
Figure~\ref{fig:sikeepercomposites} the composites that block supports. For the reduced sample,
Figure~\ref{fig:sikeepersreduced} gives the clusters and the two surviving composites,
Figures~\ref{fig:sikeeperstopping} and \ref{fig:sikeepersextra} the shot stopping views that
remain, and Figure~\ref{fig:sikeeperdist} the correlation structure and composite
distributions.

\begin{figure}[!p]
  \centering
  \includegraphics[width=\textwidth]{archive_keeper_shot_stopping.pdf}
  \caption{Saves above expectation on the primary sample, as post-shot expected goals minus
  goals allowed, which conditions on the difficulty of the shots faced rather than only their
  volume. \textbf{a}, the distribution over keeper-seasons. \textbf{b}, the best and worst
  keeper-seasons in the sample.}
  \label{fig:sikeepersfull}
\end{figure}

\begin{figure}[!p]
  \centering
  \includegraphics[width=\textwidth]{archive_keeper_corr.pdf}
  \caption{Correlation structure of the \KeeperFeatureCountArchive{} goalkeeper features on the
  primary sample, which is what allows the composite axes the reduced sample cannot build.}
  \label{fig:sikeepercorr}
\end{figure}

\begin{figure}[!p]
  \centering
  \includegraphics[width=\textwidth]{archive_keeper_composites.pdf}
  \caption{Distribution of the goalkeeper composite axes on the primary sample.}
  \label{fig:sikeepercomposites}
\end{figure}

\begin{figure}[!p]
  \centering
  \includegraphics[width=\textwidth]{keeper_clusters_pca.pdf}\par\bigskip
  \includegraphics[width=\textwidth]{keeper_composites.pdf}
  \caption{Goalkeepers on the reduced sample. Top, clusters in principal component space. Bottom,
  the two surviving composite axes.}
  \label{fig:sikeepersreduced}
\end{figure}

\begin{figure}[!p]
  \centering
  \includegraphics[width=\textwidth]{keeper_saves_vs_sota.pdf}\par\bigskip
  \includegraphics[width=\textwidth]{keeper_savepct_vs_workload.pdf}
  \caption{Shot stopping on the reduced sample. Top, saves against shots on target faced.
  Bottom, save percentage against workload. With expected goals on target withdrawn these are
  the only shot stopping views the reduced sample supports, and neither separates difficulty
  from quantity.}
  \label{fig:sikeeperstopping}
\end{figure}

\begin{figure}[!p]
  \centering
  \includegraphics[width=\textwidth]{keeper_shot_stopping_ranking.pdf}
  \caption{The shot stopping ranking that the two surviving reduced-sample composites support.}
  \label{fig:sikeepersextra}
\end{figure}

\begin{figure}[!htbp]
  \centering
  \includegraphics[width=\textwidth]{keeper_corr_heatmap.pdf}\par\bigskip
  \includegraphics[width=\textwidth]{keeper_composite_distributions.pdf}
  \caption{Reduced-sample goalkeeper features. Top, the correlation structure of what survived
  the withdrawal, which is what constrains the number of composites that can be built. Bottom,
  the distribution of the two composites.}
  \label{fig:sikeeperdist}
\end{figure}

\section{Cluster membership}
\label{sec:si:membership}

Every player's archetype assignment is published as CSV in the accompanying repository rather
than reproduced here, for both samples. The file \texttt{results/membership\_primary.csv}
carries the \ArcEligible{} eligible player-seasons of the primary sample with club, league,
season, listed position, minutes and assigned archetype, and
\texttt{results/membership\_reduced.csv} carries the reduced sample with the same fields plus
each player's distance from his cluster centroid. Both are regenerated by the pipeline, so they
cannot drift from the numbers reported above.
